\chapter{Non-Linear Supervised Alignment Methods}
\label{ch:nonlinear}

The \gls{cca} + Affine baseline from Chapter~\ref{ch:baseline} establishes a strong supervised linear upper bound, but it remains fundamentally limited to learning affine transformations between the two modalities. This chapter investigates whether relaxing that linearity constraint improves alignment. We evaluate two non-linear extensions that operate within the same supervised paradigm --- utilising the identical set of $30{,}000$ ground-truth anchor pairs --- but exploit more expressive, non-linear geometric structures to determine if they can close the remaining performance gap.

Both methods are evaluated under identical experimental conditions to those of Chapter~\ref{ch:baseline}: $n_A = 30{,}000$ anchor pairs, $k = 50$ \gls{cca} subspace, 500 held-out query pairs, averaged over three independent random seeds.

%% ─────────────────────────────────────────────────────────────────────────────
\section{Kernel CCA with Nyström Approximation}
\label{sec:kcca}

\subsection{Motivation}
\label{subsec:kcca_motivation}

Standard \gls{cca} exclusively discovers linear correlations between two spaces. However, the relationship between 3D geometric representations and natural language embeddings may contain complex non-linear structures that a purely linear projection cannot capture. Kernel CCA (KCCA) extends \gls{cca} by implicitly mapping features into a higher-dimensional Reproducing Kernel Hilbert Space (RKHS) via a kernel function. This enables the discovery of non-linear correlations without explicitly computing the high-dimensional feature map.

\subsection{Implementation}
\label{subsec:kcca_implementation}

\gls{kcca} replaces the dot product in the standard covariance computation with a kernel evaluation:
\begin{equation}
    K(x, x') = \exp\!\left(-\gamma \|x - x'\|^2\right)
    \label{eq:rbf_kernel}
\end{equation}
where $\gamma = 0.1$ is the \gls{rbf} bandwidth parameter. The resulting kernel matrix $K \in \mathbb{R}^{n \times n}$ captures non-linear similarity between samples in the original feature space.

Computing \gls{kcca} exactly requires inverting an $n \times n$ matrix, an operation that scales as $\mathcal{O}(n^3)$ and is computationally intractable for $n = 30{,}000$ anchors. To circumvent this limitation, we employ the \textbf{Nyström approximation} \cite{williams2001nystrom}, which approximates the full kernel matrix using a randomly selected subset of $m = 500$ landmark points:
\begin{equation}
    K \approx K_{nm}\, K_{mm}^{-1}\, K_{mn}
    \label{eq:nystrom}
\end{equation}
where $K_{nm} \in \mathbb{R}^{n \times m}$ contains kernel evaluations between all $n$ training points and the $m$ landmarks, and $K_{mm} \in \mathbb{R}^{m \times m}$ is the kernel matrix restricted to the landmarks themselves. This reduces the computational complexity to $\mathcal{O}(m^2 n)$, making \gls{kcca} tractable even at full dataset scale.

\subsection{Results}
\label{subsec:kcca_results}

\begin{table}[h]
\centering
\resizebox{\textwidth}{!}{%
\begin{tabular}{lcccccc}
\toprule
\textbf{Pair} & \textbf{CA Match} & \textbf{CA Top-5}
    & \textbf{KCCA Match} & \textbf{KCCA Top-5}
    & \textbf{RML Match} & \textbf{RML Top-5} \\
\midrule
PointNet $\times$ CLIP      & 13.7\% & 34.1\% & 4.7\% & 15.8\% & 9.4\%  & \textbf{27.2\%} \\
PointNet $\times$ RoBERTa   &  9.1\% & 27.8\% & 3.1\% & 13.1\% & 4.6\%  & 17.3\% \\
PointNet $\times$ BERT      &  6.9\% & 25.4\% & 3.1\% & 12.4\% & 5.1\%  & 16.5\% \\
SparseConv $\times$ CLIP    &  3.5\% & 13.2\% & 1.2\% &  6.3\% & 3.7\%  & 10.7\% \\
SparseConv $\times$ RoBERTa &  3.5\% & 13.9\% & 1.5\% &  6.1\% & 2.5\%  &  9.7\% \\
SparseConv $\times$ BERT    &  2.0\% & 10.2\% & 1.4\% &  5.4\% & 2.7\%  &  8.3\% \\
\bottomrule
\end{tabular}%
}
\caption{Non-linear methods compared to CCA+Affine (CA) across all encoder
pairs. KCCA: Kernel CCA with Nyström approximation ($\gamma$=0.1, $m$=500
landmarks). RML: Riemannian Metric Learning (batch=256, lr=0.01, 200
iterations). RML reaches 27.2\% Top-5 on PointNet$\times$CLIP, the closest
alternative to CCA+Affine at 76\% of its performance.
CA results reflect the specific CCA+Affine run executed alongside this
experiment; see the note on run-to-run variance in
Section~\ref{sec:baseline_results} (Chapter~\ref{ch:baseline}) for the
representative $35.7\%$ value used elsewhere in this thesis. Protocol:
30,000 anchor pairs, 500 query points, $k=50$, averaged over 3 random
seeds.}
\label{tab:nonlinear_results}
\end{table}

Kernel \gls{cca} with Nyström approximation achieves $15.8\%$ Top-5 accuracy on the PointNet$\times$\gls{clip} pairing. While this is a notable improvement over all unsupervised \gls{gw} methods, it remains significantly below the linear \gls{cca}+Affine baseline. This performance gap suggests that the non-linear structure captured by the \gls{rbf} kernel does not yield substantially more discriminative projections than a standard linear \gls{cca} in this specific setting. A plausible explanation is that the Nyström approximation, utilising only $m=500$ landmarks, may simply undersample the structural diversity of the full $30{,}000$-anchor distribution, thereby capping the expressivity of the resulting non-linear mapping.

%% ─────────────────────────────────────────────────────────────────────────────
\section{Riemannian Metric Learning}
\label{sec:rml}

\subsection{Motivation}
\label{subsec:rml_motivation}

Riemannian Metric Learning (RML) adopts a fundamentally different approach compared to learning an explicit mapping between the two domains. Instead, it learns a \emph{distance metric} that pulls paired anchor features as close together as possible within a shared geometric space. This is conceptually distinct from CCA+Affine, which projects both spaces into a subspace and then learns a translation between them. RML bypasses the translation step entirely, instead deforming the measurement geometry of the feature space itself so that paired 3D and text samples become intrinsically proximate.

\subsection{Implementation}
\label{subsec:rml_implementation}

RML learns a Mahalanobis distance matrix $M = L^\top L$ (symmetric positive definite, SPD) that minimises the pairwise distances between anchor pairs:
\begin{equation}
    \mathcal{L}(M) = \sum_{i=1}^{n_A} \|L(x_i - y_i)\|^2
    \label{eq:rml_loss}
\end{equation}
where $x_i$ and $y_i$ are the \gls{cca}-projected 3D and text features of the $i$-th anchor pair, and $L$ is the Cholesky factor of $M$, which guarantees positive definiteness is maintained throughout optimisation.

The matrix $L$ is optimised via stochastic gradient descent directly on the Cholesky factor, with a batch size of 256, a learning rate of 0.01, and 200 iterations. During retrieval, cosine similarity is replaced by the learned Mahalanobis distance:
\begin{equation}
    \text{sim}_M(x, y) = -\|L(x - y)\|^2
    \label{eq:mahalanobis_sim}
\end{equation}
This allows RML to learn an anisotropic geometry that weights different dimensions of the \gls{cca} subspace according to how discriminative they actually are for cross-modal retrieval.

\subsection{Results}
\label{subsec:rml_results}

RML achieves $27.2\%$ Top-5 retrieval accuracy on the PointNet$\times$\gls{clip} pair (Table~\ref{tab:nonlinear_results}). This represents the strongest result among all non-linear alternatives, approaching the \gls{cca}+Affine baseline ($35.7\%$) by capturing $76\%$ of its performance. Across every encoder combination, RML consistently outperforms \gls{kcca} by a substantial margin. This indicates that directly minimising anchor-pair distances within a learned metric space is more effective for this task than discovering non-linear correlations via kernelisation.

Figures~\ref{fig:nonlinear_pointnet} and~\ref{fig:nonlinear_sparseconv}
illustrate this comparison across all six encoder pairs, stratified by the
3D backbone as in Chapter~\ref{ch:baseline}. Both figures exhibit the exact
same hierarchical ordering regardless of the specific text encoder or 3D
backbone involved: the CCA+Affine baseline consistently achieves the highest
accuracy, KCCA the lowest, and RML reliably occupies the intermediate position.
This consistent three-way ranking across all six permutations provides robust
evidence that transcends any single metric. It confirms that the advantage of
RML over KCCA is not an artefact of one particular encoder pairing, but rather
a structural property of how the two methods leverage supervision --- one that
holds true across three distinct text encoders and two highly divergent 3D
backbones.

\begin{figure}[htpb]
  \centering
  \includegraphics[width=0.75\textwidth]{fig_nonlinear_pointnet.pdf}
  \caption{Top-5 retrieval accuracy for CCA+Affine, KCCA Nyström, and RML
  on the pre-trained PointNet backbone, across all three text encoders.
  RML consistently outperforms KCCA and approaches CCA+Affine, though
  never quite reaches it.}
  \label{fig:nonlinear_pointnet}
\end{figure}

\begin{figure}[htpb]
  \centering
  \includegraphics[width=0.75\textwidth]{fig_nonlinear_sparseconv.pdf}
  \caption{Top-5 retrieval accuracy for CCA+Affine, KCCA Nyström, and RML
  on the untrained, random-weight SparseConv backbone, across all three
  text encoders. The same three-way ranking observed in
  Figure~\ref{fig:nonlinear_pointnet} persists here despite the much
  lower absolute accuracy, reinforcing that the ordering reflects a
  genuine property of the methods rather than an artefact of one specific
  backbone.}
  \label{fig:nonlinear_sparseconv}
\end{figure}

Comparing the two figures side by side reveals a further insight. As observed in Chapter~\ref{ch:baseline}, transitioning from the pre-trained PointNet to the randomly initialised SparseConv substantially degrades absolute performance across the board: CCA+Affine's Top-5 accuracy on CLIP drops from $34.1\%$ to $13.2\%$, KCCA's from $15.8\%$ to $6.3\%$, and RML's from $27.2\%$ to $10.7\%$. However, the \emph{relative} performance gaps between the three methods remain roughly proportional. This demonstrates that the comparative advantage of RML over KCCA does not depend on a pre-trained geometric backbone; the advantage remains clearly evident even when the 3D encoder generates essentially structured noise.

The performance gap between RML ($27.2\%$) and \gls{cca}+Affine ($35.7\%$) likely stems from a fundamental difference in how the two methods utilise supervision. The affine map learns a single global linear transformation that minimises reconstruction error across \emph{all} anchor pairs simultaneously, producing one coherent global alignment. RML, by contrast, minimises local pairwise distances, an approach that may simply not generalise as robustly to unseen query pairs, despite fitting the anchor set itself quite effectively.

%% ─────────────────────────────────────────────────────────────────────────────
\section{Complete Method Ranking}
\label{sec:nonlinear_ranking}

Table~\ref{tab:complete_ranking} consolidates all alignment methods tested on PointNet$\times$\gls{clip}, providing a direct comparison across the three paradigms explored in this thesis: unsupervised GW-based methods (Chapter~\ref{ch:gw_alignment}), the supervised linear baseline (Chapter~\ref{ch:baseline}), and the supervised non-linear extensions covered in this chapter.

\begin{table}[htpb]
\centering
\renewcommand{\arraystretch}{1.2}
\resizebox{\textwidth}{!}{%
\begin{tabular}{llcccc}
\toprule
\textbf{Category} & \textbf{Method}
    & \textbf{Match} & \textbf{Top-1} & \textbf{Top-5} & \textbf{Top-10} \\
\midrule
\multicolumn{6}{l}{\textit{Unsupervised Optimal Transport}} \\
& Sliced GW ($k$=20, 5k anc.)        & 1.5\% & 1.1\% & 3.2\%           & 4.8\%  \\
& FGW ($\alpha$=0.2, 1k anc.)        & 3.6\% & 3.6\% & 4.4\%           & 5.4\%  \\
\midrule
\multicolumn{6}{l}{\textit{Supervised Linear}} \\
& CCA + Affine (30k anc.)            & 13.5\% & 14.6\% & \textbf{35.7\%} & 46.4\% \\
\midrule
\multicolumn{6}{l}{\textit{Supervised Non-Linear (this chapter)}} \\
& KCCA Nystr\"{o}m ($\gamma$=0.1)    & 4.7\%  & 4.9\%  & 15.8\%          & 23.3\% \\
& Riemannian Metric Learning         & 9.4\%  & 8.7\%  & 27.2\%          & 38.4\% \\
\midrule
\multicolumn{6}{l}{\textit{Random Baseline}} \\
& Independent Gaussian               & $\approx$0.6\% & $\approx$0.6\%
                                     & $\approx$1.6\% & $\approx$2.0\% \\
\bottomrule
\end{tabular}%
}
\caption{Complete ranking of all tested alignment methods on PointNet
$\times$ CLIP. Final Top-5 order: CCA+Affine (35.7\%) $>$ RML (27.2\%)
$>$ KCCA (15.8\%) $>$ FGW (4.4\%) $>$ SGW (3.2\%) $>$ random
($\approx$1\%). CCA+Affine result taken from
Table~\ref{tab:linear_baseline_results} (Chapter~\ref{ch:baseline}).
All results: 30,000 anchor pairs, 500 query points, $k=50$,
averaged over 3 seeds.}
\label{tab:complete_ranking}
\end{table}

The empirical pattern is unequivocal: supervised methods consistently and substantially outperform their unsupervised counterparts. Furthermore, within the supervised paradigm, the simpler linear approach remains superior to both non-linear alternatives. This finding is somewhat counter-intuitive, demonstrating that merely increasing model expressivity does not bridge the semantic gap. Rather than the sophistication of the alignment algorithm, the true performance ceiling in this problem is dictated by the availability of explicit paired supervision and the underlying geometric compatibility of the initial latent spaces.